% Abstract
\newpage
\phantomsection
\addcontentsline{toc}{matter}{ABSTRAK}
\begin{spacing}{1.15}
\begin{center}
    \textbf{ABSTRAK}
\end{center}

Model kecerdasan buatan (AI) seperti DeepBach dan Coconet dapat mengharmonisasi melodi sopran menjadi empat suara. Analisis Bach Doodle oleh Huang et al.\ (2019) memperlihatkan kuint dan oktaf sejajar pada keluaran Coconet, sedangkan buku harmoni Gustav Strube memberikan pedoman gerak suara beserta syarat dan pengecualiannya. Penelitian ini dirancang sebagai studi kasus evaluatif pada satu melodi dari bab harmonisasi \textit{chorale} Strube, dengan MODEL no.~285 sebagai calon utama setelah transkripsi lengkap diverifikasi. Sopran yang sama dipertahankan pada kedua sistem; harmonisasi cetak menjadi acuan pilihan gerak suara. Direncanakan 20 percobaan generasi per model, dengan kegagalan dan keluaran identik tetap dicatat. Program berbasis music21 membantu menemukan kandidat kejadian, kemudian peneliti membaca konteks partitur untuk membedakan penyimpangan terkonfirmasi, kejadian yang diterima menurut syarat atau pengecualian, dan kejadian yang belum dapat diputuskan. Analisis mencakup gerak melodis alto, tenor, dan bas, hubungan antarsuara, serta pola interval tersembunyi sebagai pendamping. Profil per kriteria, lokasi, dan variasi keluaran dibandingkan secara deskriptif. Penelitian diarahkan untuk menjelaskan bagian harmonisasi AI yang perlu diperiksa dalam konteks pedoman Strube. Kesimpulan dibatasi pada kasus, \textit{checkpoint}, dan pengaturan yang diuji, tanpa pengukuran persepsi pendengar atau hasil belajar pengguna.\par

\noindent\textbf{Kata Kunci:} Harmonisasi Otomatis, Kaidah Gerak Suara, Gustav Strube, \textit{Chorale}, DeepBach, Coconet
\end{spacing}
